\chapter{Sample and Apparatus}\label{chap:apparatus}
\begin{jointwork}
    The content of this chapter is adapted from Ref.\,\cite{Walter2024_LFthermal}.
\end{jointwork}

The CASPEr-Gradient-Low-Field apparatus converts a nuclear-spin signal into a detectable electrical output through three functional blocks: the magnet system sets the operating bias field and thereby the Larmor frequency of the sample; the sample and pickup coil generate and capture the magnetic response; and the SQUID and lock-in amplifier chain convert, amplify, and record the signal. This chapter describes each block in turn and closes with a summary of the key parameters.

\section{Magnet system}\label{sec:apparatus:magnet}

The apparatus is built around a flow-through liquid-helium cryostat housing a superconducting solenoid that can generate a bias field $B_0$ up to \SI{0.1}{\tesla}. The DM mass range accessible to the experiment depends on the gyromagnetic ratio of the sample and on $B_0$; for proton spins, the covered Larmor frequency range is approximately \SI{1}{\kHz} to \SI{4.3}{\MHz}.

Eight superconducting shim coils surrounding the solenoid are used to reduce the field inhomogeneity to less than \SI{2}{\ppm} over a spherical volume with radius \SI{4}{\mm}. In the measurements reported here, an inhomogeneity of approximately \SI{10}{\ppm} was achieved after tuning the shim fields using a Bayesian-optimization algorithm \cite{ShimmingPaper}. The larger-than-design inhomogeneity results from using a longer sample (\SI{24}{\mm} height) and from the insertion of an additional superconducting pickup coil.

Helmholtz coils wound outside the cryostat provide a controllable offset on $B_0$, enabling fine-tuning of the Larmor frequency without ramping the main solenoid. This arrangement is used to step the Larmor frequency in \SI{6}{\Hz} increments during the scan.

\subsection{Cryostat and superconducting magnet}\label{sec:apparatus:cryostat}

The cryostat housing the superconducting magnet is placed inside a Faraday room, shielded from external electric fields. It consists of an outer shell, reservoirs for liquid nitrogen (L\ch{N2}) and liquid helium (L\ch{He}), and vacuum insulation layers in between. Liquid nitrogen cools the cryostat and reduces the evaporation of liquid helium, while the superconducting devices and niobium shield are submerged in the L\ch{He} bath at \SI{4.2}{\kelvin}. When filled, the cryostat can maintain a superconducting environment for approximately one week. The niobium shield repels external magnetic fields, isolating internal devices from external magnetic noise.

The magnet system consists of the main solenoid coil and eight shim coils (named Z1, Z2, X, Y, XY, XZ, YZ, and X$^2$-Y$^2$), each capable of generating magnetic field components for field homogenization. The shim coil channel names indicate the geometry of the magnetic field they generate, enabling systematic correction of field inhomogeneities.

\begin{figure}[htb]
    \centering
    \includegraphics[width=0.8\textwidth]{works/master-thesis/figures/02_Apparatus/Cryostat_and_SCdevices_20220903.png}
    \caption{(a) Cryostat for maintaining superconducting temperature. L\ch{He} bath provides \SI{4.2}{\kelvin} environment for superconducting apparatus. (b) Superconducting apparatus including main coil, shim and excitation coils, SQUIDs with niobium capsule, niobium wires for connections, and gradiometer made of niobium wires.}
    \label{fig:cryostat_devices}
\end{figure}

\subsection{Excitation generation}\label{sec:apparatus:excitation}

To perform pulsed-NMR schemes, an excitation coil made of niobium wires is positioned near the sample to avoid Johnson--Nyquist noise from thermal agitation in normal-conducting materials. The excitation coil provides $\mathrm{B_x}$, $\mathrm{B_y}$, $\mathrm{B_z}$, and $\mathrm{dB_z/dz}$ channels, with the first three capable of producing homogeneous fields along the respective axes and the last generating a linear gradient along the z-axis.

Radiofrequency pulses are generated by a function generator and amplified by an rf amplifier. A tuning circuit composed of capacitors and inductors is designed for impedance matching between the function generator and the excitation coil across different frequencies.

\section{Sample}\label{sec:apparatus:sample}

Liquid methanol (\ch{CH3OH}) was chosen as the sample because of its high proton density (\SI{99}{\mmol/\cm^3}) and low freezing point (\SI{175.6}{\kelvin}), which allows it to remain liquid at the operating temperature of approximately \SI{180}{\kelvin}. With four protons per molecule, the \SI{1.2}{\cm^3} cylindrical sample (radius \SI{4}{\mm}, height \SI{24}{\mm}) contains $\num{7.2e22}$ protons.

A variable-temperature insert (VTI) positions the sample in the homogeneous region of the bias field and thermally isolates it from the liquid-helium bath. The sample temperature is regulated by flowing warm nitrogen gas through a spiral surrounding the sample holder. Temperature stability is monitored by repeated pulsed-NMR measurements: a constant NMR linewidth indicates a stable temperature.

At thermal equilibrium in the bias field $B_0$, the proton spins acquire a net longitudinal magnetization
\begin{equation}
    M_0 = \frac{n\,\gamma^2\hbar^2\,I(I+1)\,B_0}{3\,k_B\,T}
    \,,\label{eq:M0}
\end{equation}
where $n$ is the proton number density, $\gamma = 2\pi\times\SI{42.577}{\MHz/\tesla}$ is the proton gyromagnetic ratio, $I=1/2$, $k_B$ is the Boltzmann constant, and $T$ is the sample temperature. At $B_0 = \SI{31.67}{\milli\tesla}$ and $T \approx \SI{180}{\kelvin}$, the thermal polarization is $\approx\num{1.8e-7}$ and $M_0 \approx \SI{1.86e-10}{\tesla}/\mu_0$.


\section{Pickup coil and shielding}\label{sec:apparatus:pickup}

A niobium (Nb) pickup coil mounted on the outside of the VTI in the liquid-helium bath detects the transverse magnetization of the sample. The coil consists of three saddle-pair coils in a gradiometer configuration: the area of the central coil is twice that of the two outer coils, and the coils are wound in opposing senses so that far-field flux contributions cancel. The filling factor---the fraction of the coil volume occupied by the sample---is approximately \SI{8}{\%}.

A superconducting Nb shield built into the cryostat attenuates low-frequency magnetic noise. Additional shielding is provided by a $\mu$-metal shield outside the cryogenic reservoir and by a Faraday cage enclosing the entire apparatus, which suppresses electromagnetic interference (EMI). The pickup coil geometry and pulsed-NMR scheme are illustrated in Figure~\ref{fig:setup}.

\begin{figure}[htb]
    \centering
    \includegraphics[width=0.8\textwidth]{works/master-thesis/figures/02_Apparatus/Gradiometer_and_Sample_20220903.png}
    \caption{(a) Simplified schematic of gradiometer and sample positioned at the center. (b) Details of gradiometer loop connection showing the three-coil saddle-pair configuration with opposing windings for noise rejection.}
    \label{fig:gradiometer_sample}
\end{figure}

\begin{figure}[htb]
    \centering
    \includegraphics[width=\linewidth]{works/LF-thermal/figures/setup.png}
    \caption{Left: the methanol sample (blue) sits at the center of the three-coil Nb saddle-pair pickup (red). Helmholtz coils (black) fine-tune the bias field. Right: pulsed-NMR scheme used for sample characterization. The net magnetization $\mathbf{M}$ is tipped into the transverse plane by a $\pi/2$ pulse and precesses, generating a signal in the pickup and SQUID. Reproduced from Ref.\,\cite{Walter2024_LFthermal}.}
    \label{fig:setup}
\end{figure}


\section{Detection chain}\label{sec:apparatus:detection}

\subsection{SQUID detector and flux-locked-loop operation}\label{sec:apparatus:squid}

Magnetic flux in the pickup coil is inductively coupled into a direct-current superconducting quantum interference device (dc-SQUID) operated in flux-locked-loop (FLL) mode. The SQUID converts magnetic flux into electrical signals through the combination of flux quantization and Josephson tunneling effects, offering exceptional sensitivity for weak magnetic signals.

The relationship between the magnetic flux coupled to the SQUID and the feedback voltage is given by equation~\eqref{eq:squidpickup}, where the conversion factor depends on the coupling strengths and feedback parameters summarized in Table~\ref{tab:SQUIDtab}. The SQUID circuit is enclosed in a niobium capsule for shielding and protection from external disturbances.

\begin{figure}[htb]
    \centering
    \includegraphics[width=\textwidth]{works/master-thesis/figures/02_Apparatus/SQUID_Schematic_VPhi_FLL.png}
    \caption{(a) $V$--$\Phi$ characteristic of dc SQUID showing the operating point W where the feedback voltage has approximately linear dependence on magnetic flux. (b) Direct-coupled flux-locked-loop (FLL) circuit with built-in pre-amplifier and integrator for maintaining the SQUID at the linear operating point.}
    \label{fig:squid_schematic}
\end{figure}

\subsection{Signal chain and processing}\label{sec:apparatus:signal_chain}

Magnetic flux in the pickup coil is inductively coupled into a direct-current superconducting quantum interference device (dc-SQUID) operated in flux-locked-loop (FLL) mode. The properties of the SQUID sensor are listed in Table~\ref{tab:SQUIDtab}. The feedback voltage $V_\mathrm{f}$ is related to the magnetic flux in the pickup coil $\Phi_\mathrm{pick}$ by
\begin{equation}
    V_\mathrm{f}
    = \frac{R_\mathrm{f}}{M_\mathrm{f}}\,\frac{M_\mathrm{in}}{L_\mathrm{pick}+L_\mathrm{in}}\,\Phi_\mathrm{pick}
    = \varepsilon\,\Phi_\mathrm{pick}
    \,,\label{eq:squidpickup}
\end{equation}
where $R_\mathrm{f}$ is the feedback resistance, $M_\mathrm{f}$ is the feedback-coil mutual inductance, $M_\mathrm{in}$ is the mutual inductance between pickup and SQUID input coils, and $L_\mathrm{pick}$, $L_\mathrm{in}$ are the respective coil inductances. The conversion factor $\varepsilon \approx \SI{0.549}{\milli\volt/\Phi_0}$, where $\Phi_0$ is the magnetic flux quantum.

\begin{table}[ht]
    \centering
    \caption{Properties of the dc-SQUID sensor used in the measurements.}
    \begin{tabular}{lc}
        \hline\hline
        Pickup coil inductance ($L_\mathrm{pick}$) & \SI{553}{\nH} \\[0.5ex]
        SQUID input coil inductance ($L_\mathrm{in}$) & \SI{400}{\nH} \\[0.5ex]
        Mutual inductance ($M_\mathrm{in}$) & \SI{3.98}{\nH} \\[0.5ex]
        Feedback-coil inductance ($M_\mathrm{f}$) & \SI{22.83}{\Phi_0/\mA} \\[0.5ex]
        Feedback resistance ($R_\mathrm{f}$) & \SI{3}{\kohm} \\[0.5ex]
        \hline\hline
    \end{tabular}
    \label{tab:SQUIDtab}
\end{table}

The detection chain is characterized by the transfer coefficient $\beta$, defined as the ratio of the SQUID feedback-voltage amplitude $V_\perp$ to the transverse sample magnetization $\mu_0 M_\perp$:
\begin{equation}
    \beta = \frac{V_\perp}{\mu_0\,M_\perp}
    \,.\label{eq:beta}
\end{equation}
The transfer coefficient was determined at each scan step from the feedback voltage following a $\pi/2$ pulse, for which $M_\perp = M_0$. The average value is $\beta \approx 4\times10^6\,\mathrm{V/T}$.

During ALP-search measurements the SQUID and FLL electronics are battery-powered to avoid electromagnetic interference from the mains. All other electronics (lock-in amplifier, field-ramp supply) are outside the Faraday cage; only shielded BNC cables penetrating the cage are used to carry signals.

The FLL output is fed into a Zurich Instruments MFLI Lock-in Amplifier (LIA), which demodulates the signal into in-phase (I) and quadrature (Q) channels. The complex output $I + iQ$ is recorded at a sampling frequency of \SI{13.4}{\kHz} and is used to generate power spectral densities (PSDs) via fast Fourier transformation (FFT).


\section{Experimental control and data acquisition}\label{sec:apparatus:control}

Comprehensive centralized control of the CASPEr-Gradient-LF apparatus is implemented in Python to coordinate the complex timing requirements of pulsed-NMR experiments and manage communication between diverse hardware components. This control system enables synchronized operation of shim coils, rf pulse generation, and SQUID data acquisition. 

The main functions of the control system include:

\begin{enumerate}
    \item \textbf{Magnet and shim control:}
    \begin{enumerate}
        \item Remote switching of shim power supply on and off
        \item Connection and disconnection of communication between power supply and shim coils
        \item Real-time setting of shim currents to desired values, with optional heater control for current stabilization
    \end{enumerate}
    \item \textbf{Pulse generation and timing:}
    \begin{enumerate}
        \item Command interface for injecting radiofrequency pulses
        \item Recording rf pulse timing via oscilloscope into HDF5 format
        \item Analysis of function generator data and pulse characterization
    \end{enumerate}
    \item \textbf{Data acquisition and processing:}
    \begin{enumerate}
        \item Recording SQUID data and automated creation of HDF5 files for storage
        \item Real-time analysis and monitoring of acquired signals
        \item Generation and analysis of power spectral densities
    \end{enumerate}
\end{enumerate}

The control system orchestrates multiple measurement schemes including pulsed-NMR diagnostics for sample characterization, automated shimming procedures with feedback optimization, axion wind measurements, and magnetic field monitoring. Additionally, the system includes capabilities for NMR kinetic simulations, stochastic axion wind simulations, apparatus properties computation, and synchronized file management across multiple terminals.


\section{Parameters summary}\label{sec:apparatus:params}

Table~\ref{tab:parameters} summarizes the key parameters of the CASPEr-Gradient-LF apparatus for the methanol measurement reported in this thesis and for two projected future configurations using hyperpolarized samples. The thermal polarization of methanol is the dominant limiting factor in the current measurement. Hyperpolarized methanol or liquid \ch{^{129}Xe} could boost the polarization by several orders of magnitude, substantially improving the sensitivity.

\begin{table*}[ht]
\begin{threeparttable}
\centering
\caption{Parameters of the CASPEr-Gradient-LF apparatus for the current measurement (thermally polarized methanol) and two projected future configurations. Polarization is given as molar polarization (polarization$\times$concentration).}
\label{tab:parameters}
\setlength{\tabcolsep}{4pt}
\begin{tabular}{p{0.24\linewidth}
                >{\centering\arraybackslash}p{0.21\linewidth}
                >{\centering\arraybackslash}p{0.21\linewidth}
                >{\centering\arraybackslash}p{0.21\linewidth}}
\hline\hline
& This work & \multicolumn{2}{c}{Examples of future CASPEr-Gradient-LF} \\
\hline
Sample & Methanol & Methanol\tnote{a} & Liquid \ch{^{129}Xe}\tnote{b} \\
Spin density (\unit{\cm^{-3}}) & \num{5.9e22} & \num{5.9e22} & $\sim\num{1e22}$ \\
Sample volume (\unit{\cm^3}) & 1.2 & 1.2 & $\sim1$ \\
Polarization technique & Thermal & BF\tnote{c} / PHIP\tnote{d} / DNP\tnote{e} & SEOP\tnote{f} \\
Polarization & \num{1.8e-7} & $>\num{1.1e-5}$ & $\sim1$ \\
$T_2$ (\unit{\s}) & 0.72 & 1 & 1000 \\
$B_0$ inhomogeneity (\unit{\ppm}) & 10 & $\leq2$ & $\leq2$ \\
Electronic noise (\unit{\micro\Phi_0/\Hz^{1/2}}) & 7.1 & 1.0 & 1.0 \\
Flux of SPN\tnote{g} (\unit{\micro\Phi_0}) & 0.39 & 0.39 & 0.16 \\
Scan range & \SI{1.348450}{\MHz}--\SI{1.348690}{\MHz} & $\sim$\SI{1}{\kHz}--\SI{4.3}{\MHz} & $\sim$\SI{1}{\kHz}--\SI{1.2}{\MHz} \\
$|g_\mathrm{aNN}|$ sensitivity (\unit{\GeV^{-1}}) & \num{3e-2} & \num{1.0e-6}--\num{1.2e-5} & \num{1e-12}--\num{1e-11} \\
\hline\hline
\end{tabular}
\begin{tablenotes}
\footnotesize
\item[a] Methanol at \SI{180}{\kelvin} prepolarized at \SI{2}{\tesla}; measurement time $100\,\tau_a$ per scan step.
\item[b] \ch{^{129}Xe}-enriched liquid sample.
\item[c] Brute-force prepolarization in a strong field at low temperature.
\item[d] Parahydrogen-induced polarization (PHIP).
\item[e] Dynamic nuclear polarization (DNP).
\item[f] Spin-exchange optical pumping (SEOP).
\item[g] Spin-projection noise (SPN).
\end{tablenotes}
\end{threeparttable}
\end{table*}
